\documentclass[11pt]{article}

\usepackage[a4paper,margin=1in]{geometry}
\usepackage{amsmath,amssymb,amsthm,mathtools}
\usepackage{bm}
\usepackage{enumitem}
\usepackage{hyperref}

\numberwithin{equation}{section}

\newtheorem{theorem}{Theorem}[section]
\newtheorem{proposition}[theorem]{Proposition}
\newtheorem{lemma}[theorem]{Lemma}
\newtheorem{conjecture}[theorem]{Conjecture}
\newtheorem{remark}[theorem]{Remark}
\newtheorem{example}[theorem]{Example}

\DeclareMathOperator{\Tr}{Tr}
\DeclareMathOperator{\KL}{KL}
\DeclareMathOperator{\Sym}{Sym}
\DeclareMathOperator{\diag}{diag}
\DeclareMathOperator{\argmin}{argmin}

\title{A Counterexample to the Fourth-Order Operator-Norm Route for Dimension-Free Local Convergence of Gaussian Natural Gradient Flow}
\author{}
\date{}

\begin{document}

\maketitle

\begin{abstract}
We analyze the local convergence rate of the Gaussian natural gradient flow near the optimal Gaussian variational approximation. 
In equilibrium-whitened coordinates, the current manuscript proves a local rate lower bound involving a fourth-order operator norm. 
A natural conjecture was that this fourth-order operator norm admits a dimension-free logarithmic bound in the condition number. 
We show that this stronger operator-norm conjecture is false by constructing a family of admissible radial potentials for which the fourth-order trace mode exceeds the proposed one-dimensional tail-cut bound. 
However, the same examples are radial and even, so the mean-covariance coupling vanishes and the actual local convergence rate remains maximal. 
Therefore the counterexample does not refute the dimension-free local-rate conjecture itself. 
Instead, it shows that controlling the largest fourth-order covariance mode is too lossy. 
The correct rate-controlling quantity should measure the mean-covariance coupling, which can be formulated using only the Hessian field.
\end{abstract}

\section{Problem setup}

We consider a target posterior distribution on $\mathbb R^{N_\theta}$ of the form
\begin{equation}
    \rho_{\mathrm{post}}(\theta) \propto \exp(-V(\theta)),
\end{equation}
where $V:\mathbb R^{N_\theta}\to \mathbb R$ is the negative log-density up to an additive constant. 
The Gaussian variational family is
\begin{equation}
    \mathcal P_G
    =
    \left\{
    \rho_a = \mathcal N(m,C):
    a=(m,C)\in \mathbb R^{N_\theta}\times \Sym_{++}(N_\theta)
    \right\}.
\end{equation}
The Gaussian variational inference problem is
\begin{equation}
    a_\star=(m_\star,C_\star)
    \in
    \argmin_a \KL(\rho_a\|\rho_{\mathrm{post}}).
\end{equation}

The Gaussian natural gradient flow, equivalently the Fisher--Rao gradient flow restricted to the Gaussian family, is
\begin{align}
    \frac{d m_t}{dt}
    &=
    C_t
    \mathbb E_{\theta\sim \mathcal N(m_t,C_t)}
    [
    \nabla_\theta \log \rho_{\mathrm{post}}(\theta)
    ],
    \\
    \frac{d C_t}{dt}
    &=
    C_t
    +
    C_t
    \mathbb E_{\theta\sim \mathcal N(m_t,C_t)}
    [
    \nabla_\theta\nabla_\theta \log \rho_{\mathrm{post}}(\theta)
    ]
    C_t.
\end{align}
Since $\log \rho_{\mathrm{post}}=-V+\mathrm{const.}$, this can equivalently be written as
\begin{align}
    \frac{d m_t}{dt}
    &=
    -
    C_t
    \mathbb E_{\theta\sim \mathcal N(m_t,C_t)}
    [
    \nabla V(\theta)
    ],
    \\
    \frac{d C_t}{dt}
    &=
    C_t
    -
    C_t
    \mathbb E_{\theta\sim \mathcal N(m_t,C_t)}
    [
    \nabla^2 V(\theta)
    ]
    C_t.
\end{align}

Throughout this report, we assume the standard log-concavity and smoothness condition
\begin{equation}
    \alpha I
    \preceq
    \nabla^2 V(\theta)
    \preceq
    \beta I,
    \qquad
    \theta\in\mathbb R^{N_\theta},
\end{equation}
with condition number
\begin{equation}
    \kappa=\frac{\beta}{\alpha}.
\end{equation}

We focus on the local convergence rate near the Gaussian variational optimum. 
By affine invariance, we may work in equilibrium-whitened coordinates
\begin{equation}
    z=C_\star^{-1/2}(\theta-m_\star),
\end{equation}
so that the stationary Gaussian becomes
\begin{equation}
    \rho_{a_\star}=\mathcal N(0,I).
\end{equation}
In these coordinates, writing again $V$ for the transformed potential, the stationarity conditions become
\begin{equation}
    \mathbb E[\nabla V(z)] = 0,
    \qquad
    \mathbb E[\nabla^2 V(z)] = I,
    \qquad
    z\sim \mathcal N(0,I).
\end{equation}
The Hessian bounds become
\begin{equation}
    \alpha I
    \preceq
    \nabla^2 V(z)
    \preceq
    \beta I,
    \qquad
    \alpha\le 1\le \beta.
\end{equation}
We denote
\begin{equation}
    W(z):=\nabla^2 V(z).
\end{equation}

The goal is to understand the largest eigenvalue $\lambda_{\star,\max}<0$ of the linearized Jacobian $J_\star$ at equilibrium. 
The local convergence rate is controlled by
\begin{equation}
    -\lambda_{\star,\max}.
\end{equation}
Larger values of $-\lambda_{\star,\max}$ correspond to faster local convergence.

\section{Linearization and the rate-controlling quadratic form}

Let
\begin{equation}
    m=u,
    \qquad
    C=I+X,
    \qquad
    X\in \Sym(N_\theta),
\end{equation}
denote a perturbation around the equilibrium $(0,I)$. 
The linearized dynamics has the form
\begin{equation}
    -J_\star(u,X)
    =
    \left(
    u+\frac12 T[X],
    \,
    X+T^\ast[u]+\frac12 H[X]
    \right),
\end{equation}
where the linear operators can be written using only the Hessian $W(z)$:
\begin{align}
    T[X]
    &=
    \mathbb E
    \left[
    \Tr(XW(z))z
    \right],
    \\
    T^\ast[u]
    &=
    \mathbb E
    \left[
    (u^\top z)W(z)
    \right],
    \\
    H[X]
    &=
    \mathbb E
    \left[
    W(z)\left(z^\top Xz-\Tr X\right)
    \right].
\end{align}
These are Hessian-only formulas. 
If $V$ is smoother, the same objects can be rewritten using third and fourth derivatives by Gaussian integration by parts, but such derivative notation is not essential for the formulation.

Define the bilinear form
\begin{equation}
    B(S,S')
    :=
    \mathbb E
    \left[
    (Sz)^\top (W(z)-I)(S'z)
    \right],
    \qquad
    S,S'\in \Sym(N_\theta).
\end{equation}
Then
\begin{equation}
    B(S,S)
    =
    \Tr(SH[S]).
\end{equation}

For diagonal $X=\diag(\lambda)$, define
\begin{equation}
    G_{ij}
    =
    \mathbb E[W_{ii}(z)z_j^2],
    \qquad
    A:=G-11^\top.
\end{equation}
Then
\begin{equation}
    \lambda^\top A\lambda
    =
    B(\diag\lambda,\diag\lambda).
\end{equation}

The fourth-order operator norm of interest is
\begin{equation}
    \Lambda
    :=
    \sup_{\|S\|_F=1}
    B(S,S).
\end{equation}
Equivalently, in a fixed diagonal frame,
\begin{equation}
    \Lambda
    \ge
    \lambda_{\max}(G-11^\top),
\end{equation}
and after optimizing over orthonormal frames, $\Lambda$ captures the largest covariance-mode curvature contribution.

The existing manuscript proof bounds the local rate through an estimate on this fourth-order operator. 
In particular, it proves a bound of the form
\begin{equation}
    -\lambda_{\star,\max}
    \ge
    \frac{1}{4+\Gamma},
\end{equation}
where
\begin{equation}
    \Gamma
    =
    \min
    \left\{
    \beta-1,
    \,
    N_\theta
    \left(
    3+\frac{4}{\sqrt{\pi}}(1+\log\kappa)
    \right)
    \right\}.
\end{equation}
The second branch contains the dimension factor $N_\theta$. 
The natural question is whether this factor is intrinsic or merely a proof artifact.

\section{The dimension-free local-rate conjecture}

The broad conjecture motivating the analysis is the following.

\begin{conjecture}[Dimension-free local rate]
Under the equilibrium-whitened assumptions
\begin{equation}
    \mathbb E[\nabla V(z)]=0,
    \qquad
    \mathbb E[\nabla^2 V(z)]=I,
    \qquad
    \alpha I\preceq \nabla^2 V(z)\preceq \beta I,
\end{equation}
there exists a constant $c(\kappa)>0$, independent of $N_\theta$, such that
\begin{equation}
    -\lambda_{\star,\max}
    \ge
    c(\kappa).
\end{equation}
The expected sharp dependence is logarithmic:
\begin{equation}
    c(\kappa)
    \gtrsim
    \frac{1}{1+\log\kappa}.
\end{equation}
\end{conjecture}

The counterexample below does not contradict this conjecture.

A stronger conjecture was proposed as a possible route to proving the dimension-free local-rate conjecture. 
This stronger conjecture attempts to control the fourth-order operator norm $\Lambda$ directly.

Define the one-dimensional tail-cut constant
\begin{equation}
    D_\kappa
    :=
    \left(
    4+\frac{4}{\sqrt{\pi}}
    \right)(1+\log\kappa).
\end{equation}
This constant arises from a one-dimensional Gaussian tail-cut estimate. 
Roughly, if
\begin{equation}
    f:\mathbb R\to[\alpha,\beta],
    \qquad
    \mathbb E[f(Z)]=1,
    \qquad
    Z\sim \mathcal N(0,1),
\end{equation}
then the tail-cut argument gives a dimension-free upper bound of the form
\begin{equation}
    \mathbb E[f(Z)Z^2]\le D_\kappa.
\end{equation}
This motivates the following operator-norm conjecture.

\begin{conjecture}[Fourth-order operator-norm conjecture]
Under the same equilibrium-whitened assumptions,
\begin{equation}
    \Lambda
    \le
    D_\kappa-1.
\end{equation}
\end{conjecture}

This conjecture would imply a dimension-free control of the fourth-order covariance operator and hence would give a route to proving a dimension-free local rate. 
However, this fourth-order operator-norm conjecture is false.

The counterexample below shows that $\Lambda$ can exceed $D_\kappa-1$, even for true Hessian fields $W=\nabla^2 V$ coming from admissible strongly convex smooth potentials. 
Therefore, the proof strategy based on bounding $\Lambda$ by $D_\kappa-1$ cannot work.

Importantly, the counterexample does not produce slow local convergence. 
In fact, the local rate remains maximal. 
Thus the example refutes the fourth-order operator-norm conjecture, but not the dimension-free local-rate conjecture.

\section{Radial potentials and their Hessian structure}

We construct a radial potential
\begin{equation}
    V(z)=g(\|z\|^2),
    \qquad
    z\in\mathbb R^{N_\theta}.
\end{equation}
Let
\begin{equation}
    s=\|z\|^2,
    \qquad
    \rho=\|z\|=\sqrt{s}.
\end{equation}

For radial potentials, the Hessian $W(z)=\nabla^2 V(z)$ has two eigenvalues:
\begin{enumerate}[label=(\roman*)]
    \item a tangential eigenvalue $h(s)$, with multiplicity $N_\theta-1$;
    \item a radial eigenvalue $r(s)$, in the direction $z/\|z\|$.
\end{enumerate}
Explicitly,
\begin{equation}
    h(s)=2g'(s),
\end{equation}
and
\begin{equation}
    r(s)=2g'(s)+4s g''(s).
\end{equation}
Equivalently,
\begin{equation}
    W(z)
    =
    h(s)I
    +
    (r(s)-h(s))
    \frac{zz^\top}{\|z\|^2},
    \qquad
    z\neq 0.
\end{equation}

The relation between $h$ and $r$ is the Hardy-average identity
\begin{equation}
    h(\rho^2)
    =
    \frac{1}{\rho}
    \int_0^\rho r(t^2)\,dt.
\end{equation}
Therefore, once the radial eigenvalue $r$ is specified, the tangential eigenvalue $h$ is determined by the Hardy average.

The Hessian bound
\begin{equation}
    \alpha I
    \preceq
    W(z)
    \preceq
    \beta I
\end{equation}
is equivalent to
\begin{equation}
    \alpha\le h(s)\le \beta,
    \qquad
    \alpha\le r(s)\le \beta.
\end{equation}
Thus, to construct an admissible radial potential, it suffices to choose $r\in[\alpha,\beta]$ such that its Hardy average $h$ also lies in $[\alpha,\beta]$.

\section{The radial threshold construction}

Fix $0<\alpha<1<\beta$, and let
\begin{equation}
    \kappa=\frac{\beta}{\alpha}.
\end{equation}
Choose a threshold $s_0>0$, and define the radial eigenvalue
\begin{equation}
    r(s)
    =
    \alpha \mathbf 1_{\{s\le s_0\}}
    +
    \beta \mathbf 1_{\{s>s_0\}}.
\end{equation}
Using the Hardy-average formula, the corresponding tangential eigenvalue is
\begin{equation}
    h(s)
    =
    \beta
    -
    (\beta-\alpha)
    \min
    \left\{
    1,\sqrt{\frac{s_0}{s}}
    \right\}.
\end{equation}
Indeed, if $s\le s_0$, then $h(s)=\alpha$. 
If $s>s_0$, then
\begin{equation}
    h(s)
    =
    \frac{1}{\sqrt{s}}
    \left(
    \int_0^{\sqrt{s_0}}\alpha\,dt
    +
    \int_{\sqrt{s_0}}^{\sqrt{s}}\beta\,dt
    \right)
    =
    \beta-(\beta-\alpha)\sqrt{\frac{s_0}{s}}.
\end{equation}
Therefore
\begin{equation}
    \alpha\le h(s)\le \beta,
    \qquad
    \alpha\le r(s)\le \beta.
\end{equation}
Thus the Hessian eigenvalues are uniformly bounded between $\alpha$ and $\beta$.

Strictly speaking, the sharp threshold profile is not smooth. 
This is not a serious obstruction. 
One may replace $r$ by a smooth monotone transition $r_\varepsilon$ from $\alpha$ to $\beta$, define $h_\varepsilon$ by the Hardy average, and then define $g_\varepsilon$ by
\begin{equation}
    g_\varepsilon'(s)=\frac12 h_\varepsilon(s).
\end{equation}
The resulting $V_\varepsilon(z)=g_\varepsilon(\|z\|^2)$ is smooth, satisfies the same Hessian bounds, and converges to the threshold construction as $\varepsilon\to 0$. 
Therefore the threshold profile should be understood as the limiting model of a smooth admissible family.

\section{Whitening condition}

We now impose the equilibrium-whitened condition
\begin{equation}
    \mathbb E[W(z)] = I,
    \qquad
    z\sim \mathcal N(0,I).
\end{equation}
For radial potentials, this condition reduces to a scalar condition. 
Let
\begin{equation}
    m=N_\theta+2,
    \qquad
    S_m\sim \chi_m^2.
\end{equation}
Then
\begin{equation}
    \mathbb E[W(z)]=I
    \quad
    \Longleftrightarrow
    \quad
    \mathbb E_{\chi_m^2}[h(S_m)] = 1.
\end{equation}

For the threshold family,
\begin{equation}
    h(s)
    =
    \beta
    -
    (\beta-\alpha)
    \min
    \left\{
    1,\sqrt{\frac{s_0}{s}}
    \right\}.
\end{equation}
Thus the whitening condition becomes
\begin{equation}
    \beta
    -
    (\beta-\alpha)
    \mathbb E_{\chi_m^2}
    \left[
    \min
    \left\{
    1,\sqrt{\frac{s_0}{S_m}}
    \right\}
    \right]
    =
    1.
\end{equation}
Equivalently,
\begin{equation}
    \mathbb E_{\chi_m^2}
    \left[
    \min
    \left\{
    1,\sqrt{\frac{s_0}{S_m}}
    \right\}
    \right]
    =
    \frac{\beta-1}{\beta-\alpha}.
\end{equation}
Define
\begin{equation}
    p:=\frac{1-\alpha}{\beta-\alpha}.
\end{equation}
Then
\begin{equation}
    \frac{\beta-1}{\beta-\alpha}
    =
    1-p.
\end{equation}

Let $G_m$ denote the cumulative distribution function of $\chi_m^2$, and define
\begin{equation}
    T_m
    :=
    \frac{\Gamma((m-1)/2)}
    {\sqrt{2}\Gamma(m/2)}.
\end{equation}
A direct computation gives
\begin{equation}
    \mathbb E_{\chi_m^2}
    \left[
    \min
    \left\{
    1,\sqrt{\frac{s_0}{S_m}}
    \right\}
    \right]
    =
    G_m(s_0)
    +
    \sqrt{s_0}T_m
    \left(1-G_{m-1}(s_0)\right).
\end{equation}
Therefore $s_0$ is chosen as the unique solution of
\begin{equation}
    \Psi(s_0)
    =
    1-p,
\end{equation}
where
\begin{equation}
    \Psi(s)
    :=
    G_m(s)
    +
    \sqrt{s}T_m
    \left(1-G_{m-1}(s)\right).
\end{equation}
Since $\Psi(0)=0$, $\Psi(\infty)=1$, and $\Psi$ is monotone increasing, such $s_0$ exists and is unique.

With this choice of $s_0$, the radial threshold potential satisfies
\begin{equation}
    \mathbb E[\nabla^2 V(z)] = I.
\end{equation}
Thus it is admissible in equilibrium-whitened coordinates.

\section{The fourth-order trace mode}

Now consider the normalized trace direction
\begin{equation}
    S_0=\frac{I}{\sqrt{N_\theta}}.
\end{equation}
Recall that
\begin{equation}
    B(S,S)
    =
    \mathbb E
    \left[
    (Sz)^\top (W(z)-I)(Sz)
    \right].
\end{equation}
For $S_0=I/\sqrt{N_\theta}$,
\begin{equation}
    B(S_0,S_0)
    =
    \frac{1}{N_\theta}
    \mathbb E
    [
    z^\top (W(z)-I)z
    ].
\end{equation}
Since $W(z)$ has radial eigenvalue $r(s)$ in the direction $z$,
\begin{equation}
    z^\top W(z)z
    =
    r(s)\|z\|^2
    =
    r(s)s.
\end{equation}
Therefore
\begin{equation}
    B(S_0,S_0)
    =
    \frac{1}{N_\theta}
    \mathbb E[(r(S)-1)S],
    \qquad
    S\sim \chi_{N_\theta}^2.
\end{equation}
By the standard size-bias identity for chi-square variables,
\begin{equation}
    \frac{1}{N_\theta}
    \mathbb E_{\chi_{N_\theta}^2}[r(S)S]
    =
    \mathbb E_{\chi_{N_\theta+2}^2}[r(S)].
\end{equation}
Hence
\begin{equation}
    B(S_0,S_0)
    =
    \mathbb E_{\chi_{N_\theta+2}^2}[r(S)]-1.
\end{equation}
For the threshold profile,
\begin{equation}
    r(s)
    =
    \alpha \mathbf 1_{\{s\le s_0\}}
    +
    \beta \mathbf 1_{\{s>s_0\}},
\end{equation}
we obtain
\begin{equation}
    B(S_0,S_0)
    =
    \alpha G_m(s_0)
    +
    \beta(1-G_m(s_0))
    -1.
\end{equation}
Equivalently,
\begin{equation}
    B(S_0,S_0)
    =
    (\beta-1)
    -
    (\beta-\alpha)G_m(s_0),
    \qquad
    m=N_\theta+2.
\end{equation}

Since
\begin{equation}
    \Lambda
    =
    \sup_{\|S\|_F=1}B(S,S),
\end{equation}
we have the lower bound
\begin{equation}
    \Lambda
    \ge
    B(S_0,S_0)
    =
    (\beta-1)
    -
    (\beta-\alpha)G_m(s_0).
\end{equation}
This is the key estimate.

\section{Asymptotic growth of the trace mode}

For fixed $\alpha,\beta$ with $0<\alpha<1<\beta$, the threshold $s_0$ is determined by
\begin{equation}
    \Psi(s_0)=1-p,
    \qquad
    p=\frac{1-\alpha}{\beta-\alpha}.
\end{equation}
As $N_\theta\to\infty$, the random variable
\begin{equation}
    S_m\sim \chi_m^2,
    \qquad
    m=N_\theta+2,
\end{equation}
concentrates around $m$. 
Since
\begin{equation}
    \min\left\{1,\sqrt{\frac{s_0}{S_m}}\right\}
\end{equation}
is approximately $\sqrt{s_0/m}$ when $s_0<m$, the whitening condition heuristically gives
\begin{equation}
    \sqrt{\frac{s_0}{m}}
    \approx
    1-p.
\end{equation}
Thus
\begin{equation}
    \frac{s_0}{m}
    \approx
    (1-p)^2
    <
    1.
\end{equation}
Therefore $s_0$ lies a fixed fraction below the typical value of $\chi_m^2$. 
Consequently,
\begin{equation}
    G_m(s_0)
    =
    \mathbb P(\chi_m^2\le s_0)
    \to 0
    \qquad
    \text{as }
    N_\theta\to\infty.
\end{equation}
It follows that
\begin{equation}
    B(S_0,S_0)
    =
    (\beta-1)
    -
    (\beta-\alpha)G_m(s_0)
    \to
    \beta-1.
\end{equation}
Therefore,
\begin{equation}
    \Lambda
    \ge
    B(S_0,S_0)
    \to
    \beta-1.
\end{equation}
If $\beta-1>D_\kappa-1$, then for all sufficiently large dimensions,
\begin{equation}
    \Lambda>D_\kappa-1.
\end{equation}
This contradicts the operator-norm conjecture
\begin{equation}
    \Lambda\le D_\kappa-1.
\end{equation}

\section{Explicit violation}

Choose, for example,
\begin{equation}
    \alpha=\frac13,
    \qquad
    \beta=\frac{\kappa}{3}.
\end{equation}
Then
\begin{equation}
    \frac{\beta}{\alpha}=\kappa.
\end{equation}

For $\kappa=10^3$,
\begin{equation}
    D_\kappa-1
    =
    \left(4+\frac{4}{\sqrt{\pi}}\right)(1+\log 10^3)-1
    \approx
    48.48.
\end{equation}
At dimension $N_\theta=1024$, the radial threshold construction gives
\begin{equation}
    B(S_0,S_0)
    \approx
    57.23.
\end{equation}
Therefore,
\begin{equation}
    \Lambda
    \ge
    B(S_0,S_0)
    \approx
    57.23
    >
    48.48
    \approx
    D_\kappa-1.
\end{equation}

For $\kappa=10^4$, again with
\begin{equation}
    \alpha=\frac13,
    \qquad
    \beta=\frac{10^4}{3},
    \qquad
    N_\theta=1024,
\end{equation}
one obtains
\begin{equation}
    D_\kappa-1
    \approx
    62.88,
\end{equation}
whereas
\begin{equation}
    B(S_0,S_0)
    \approx
    82.58.
\end{equation}
Thus
\begin{equation}
    \Lambda
    \ge
    82.58
    >
    62.88
    =
    D_\kappa-1.
\end{equation}
Hence the conjectured operator-norm bound fails.

\section{Why this does not contradict dimension-free local convergence}

The above example is radial. 
Therefore
\begin{equation}
    V(-z)=V(z).
\end{equation}
By symmetry, the mean-covariance coupling vanishes:
\begin{equation}
    T[X]
    =
    \mathbb E[
    \Tr(XW(z))z
    ]
    =
    0.
\end{equation}
Indeed, $W(z)$ is even in $z$, while $z$ is odd, so the integrand is odd.

Therefore the linearized operator decouples into the mean block and covariance block:
\begin{equation}
    -J_\star(u,X)
    =
    \left(
    u,
    \,
    X+\frac12 H[X]
    \right).
\end{equation}
The mean block has contraction rate $1$. 
The covariance block is controlled by the spectrum of $H$.

For this radial family, $H$ has two relevant isotypic eigenvalues:
\begin{enumerate}[label=(\roman*)]
    \item the trace mode,
    \begin{equation}
        B_{\mathrm{tr}}
        =
        B\left(\frac{I}{\sqrt{N_\theta}},\frac{I}{\sqrt{N_\theta}}\right)
        =
        \mathbb E_{\chi_{N_\theta+2}^2}[r]-1;
    \end{equation}
    \item the traceless covariance mode,
    \begin{equation}
        c_{\mathrm{tl}}
        =
        2
        \left(
        \mathbb E_{\chi_{N_\theta+4}^2}[h]-1
        \right).
    \end{equation}
\end{enumerate}

The trace mode is precisely the mode used to make $\Lambda$ large. 
But it is positive. 
It is therefore a fast-contracting mode, not a slow mode.

Moreover, $h$ is nondecreasing, and the whitening condition gives
\begin{equation}
    \mathbb E_{\chi_{N_\theta+2}^2}[h]=1.
\end{equation}
Since $\chi_{N_\theta+4}^2$ stochastically dominates $\chi_{N_\theta+2}^2$, we have
\begin{equation}
    \mathbb E_{\chi_{N_\theta+4}^2}[h]
    \ge
    \mathbb E_{\chi_{N_\theta+2}^2}[h]
    =
    1.
\end{equation}
Thus
\begin{equation}
    c_{\mathrm{tl}}\ge 0.
\end{equation}
Consequently,
\begin{equation}
    \lambda_{\min}(H)\ge 0.
\end{equation}
The local covariance-block rate is therefore at least $1$, and in the normalized convention,
\begin{equation}
    -\lambda_{\star,\max}
    =
    1.
\end{equation}

Thus the same example that violates
\begin{equation}
    \Lambda\le D_\kappa-1
\end{equation}
has maximal local convergence rate.

This is the central point: the counterexample does not produce slow convergence. 
Instead, it shows that $\Lambda$ is the wrong quantity to control. 
The quantity $\Lambda$ measures the largest positive covariance curvature mode, which may correspond to fast contraction. 
The slow rate is instead controlled by the bottom of the relevant linearized spectrum and by the mean-covariance coupling.

\section{Interpretation}

The operator-norm conjecture fails because it tries to control the largest eigenvalue of a fourth-order covariance operator:
\begin{equation}
    \Lambda=\sup_{\|S\|_F=1}B(S,S).
\end{equation}
But the local convergence rate is not determined by this largest eigenvalue alone. 
A large positive value of $B(S,S)$ can make the flow faster in direction $S$. 
Therefore bounding $\Lambda$ is a sufficient but structurally lossy strategy.

The radial example concentrates the large value of $B(S,S)$ on the trace/volume mode. 
This makes $\Lambda$ large, but does not slow down the flow. 
Because the potential is symmetric, the dangerous mean-covariance coupling is absent.

Therefore the correct conclusion is:
\begin{equation}
    \boxed{
    \text{The conjecture } \Lambda\le D_\kappa-1 \text{ is false.}
    }
\end{equation}
But also:
\begin{equation}
    \boxed{
    \text{The dimension-free local-rate conjecture remains open.}
    }
\end{equation}
The counterexample only rules out one proof strategy.

\section{Revised conjectural direction}

The failure of the $\Lambda$-bound suggests that the local rate should be controlled by a lower-order, coupling-sensitive quantity rather than by the fourth-order operator norm.

A Hessian-only candidate is
\begin{equation}
    \tau_H
    :=
    \sup_{\substack{|w|=1\\ \|S\|_F=1}}
    \left|
    \mathbb E[
    w^\top W(z)Sz
    ]
    \right|.
\end{equation}
Equivalently,
\begin{equation}
    \tau_H
    =
    \sup_{|w|=1}
    \left\|
    \mathbb E[
    W(z)(w^\top z)
    ]
    \right\|_F.
\end{equation}
This quantity measures the first Hermite coefficient of the Hessian field $W(z)$. 
It captures the mean-covariance coupling $T$, since
\begin{equation}
    T[X]
    =
    \mathbb E[
    \Tr(XW(z))z
    ].
\end{equation}

For radial even potentials,
\begin{equation}
    \tau_H=0,
\end{equation}
which correctly predicts that the counterexample should not slow the flow.

A more plausible dimension-free conjecture is therefore the following.

\begin{conjecture}[Hessian first-moment conjecture]
Under the equilibrium-whitened assumptions
\begin{equation}
    \mathbb E[W(z)]=I,
    \qquad
    \alpha I\preceq W(z)\preceq \beta I,
\end{equation}
there exists a universal constant $C>0$ such that
\begin{equation}
    \tau_H
    \le
    C(1-\alpha)\sqrt{1+\log\kappa},
\end{equation}
or at least
\begin{equation}
    \tau_H^2
    \le
    C(1-\alpha)^2(1+\log\kappa),
\end{equation}
dimension-free in $N_\theta$.
\end{conjecture}

Such a bound would imply a fast dimension-free local rate of the form
\begin{equation}
    -\lambda_{\star,\max}
    \gtrsim
    \frac{1}{1+(1-\alpha)^2(1+\log\kappa)}.
\end{equation}
Unlike the false $\Lambda$-conjecture, this conjecture targets the part of the linearization that can genuinely slow down convergence: the coupling between the mean and covariance blocks.

\section{Conclusion}

The radial threshold construction gives a rigorous counterexample to the fourth-order operator-norm conjecture
\begin{equation}
    \Lambda\le D_\kappa-1.
\end{equation}
The construction is admissible because its Hessian eigenvalues lie in $[\alpha,\beta]$, and the threshold $s_0$ is chosen so that the equilibrium-whitened condition
\begin{equation}
    \mathbb E[\nabla^2 V(z)]=I
\end{equation}
holds. 
In the normalized trace direction $S_0=I/\sqrt{N_\theta}$, the fourth-order form satisfies
\begin{equation}
    B(S_0,S_0)
    =
    (\beta-1)-(\beta-\alpha)G_{N_\theta+2}(s_0).
\end{equation}
As $N_\theta\to\infty$, this approaches $\beta-1$. 
For parameter regimes such as
\begin{equation}
    \alpha=\frac13,
    \qquad
    \beta=\frac{\kappa}{3},
\end{equation}
this quantity exceeds
\begin{equation}
    D_\kappa-1
    =
    \left(4+\frac{4}{\sqrt{\pi}}\right)(1+\log\kappa)-1
\end{equation}
in sufficiently high dimension. 
Hence
\begin{equation}
    \Lambda>D_\kappa-1,
\end{equation}
contradicting the proposed operator-norm conjecture.

However, the construction is radial and even. 
Therefore the mean-covariance coupling vanishes, the large trace mode is a fast mode, and the local convergence rate remains maximal:
\begin{equation}
    -\lambda_{\star,\max}=1.
\end{equation}
Thus the counterexample does not refute the dimension-free local-rate conjecture. 
It shows that the proof strategy based on bounding the largest fourth-order covariance mode is too lossy. 
The correct object should instead measure the mean-covariance coupling, for example through the Hessian first-moment quantity
\begin{equation}
    \tau_H
    =
    \sup_{\substack{|w|=1\\ \|S\|_F=1}}
    \left|
    \mathbb E[
    w^\top \nabla^2 V(z)Sz
    ]
    \right|.
\end{equation}
The remaining open problem is to prove a dimension-free logarithmic bound on this coupling quantity under only the Hessian bounds and the whitening condition.

\end{document}